% Options for packages loaded elsewhere
\PassOptionsToPackage{unicode}{hyperref}
\PassOptionsToPackage{hyphens}{url}
\documentclass[
  ignorenonframetext,
]{beamer}
\newif\ifbibliography
\usepackage{pgfpages}
\setbeamertemplate{caption}[numbered]
\setbeamertemplate{caption label separator}{: }
\setbeamercolor{caption name}{fg=normal text.fg}
\beamertemplatenavigationsymbolsempty
% remove section numbering
\setbeamertemplate{part page}{
  \centering
  \begin{beamercolorbox}[sep=16pt,center]{part title}
    \usebeamerfont{part title}\insertpart\par
  \end{beamercolorbox}
}
\setbeamertemplate{section page}{
  \centering
  \begin{beamercolorbox}[sep=12pt,center]{section title}
    \usebeamerfont{section title}\insertsection\par
  \end{beamercolorbox}
}
\setbeamertemplate{subsection page}{
  \centering
  \begin{beamercolorbox}[sep=8pt,center]{subsection title}
    \usebeamerfont{subsection title}\insertsubsection\par
  \end{beamercolorbox}
}
% Prevent slide breaks in the middle of a paragraph
\widowpenalties 1 10000
\raggedbottom
\AtBeginPart{
  \frame{\partpage}
}
\AtBeginSection{
  \ifbibliography
  \else
    \frame{\sectionpage}
  \fi
}
\AtBeginSubsection{
  \frame{\subsectionpage}
}
\usepackage{iftex}
\ifPDFTeX
  \usepackage[T1]{fontenc}
  \usepackage[utf8]{inputenc}
  \usepackage{textcomp} % provide euro and other symbols
\else % if luatex or xetex
  \usepackage{unicode-math} % this also loads fontspec
  \defaultfontfeatures{Scale=MatchLowercase}
  \defaultfontfeatures[\rmfamily]{Ligatures=TeX,Scale=1}
\fi
\usepackage{lmodern}
\ifPDFTeX\else
  % xetex/luatex font selection
\fi
% Use upquote if available, for straight quotes in verbatim environments
\IfFileExists{upquote.sty}{\usepackage{upquote}}{}
\IfFileExists{microtype.sty}{% use microtype if available
  \usepackage[]{microtype}
  \UseMicrotypeSet[protrusion]{basicmath} % disable protrusion for tt fonts
}{}
\makeatletter
\@ifundefined{KOMAClassName}{% if non-KOMA class
  \IfFileExists{parskip.sty}{%
    \usepackage{parskip}
  }{% else
    \setlength{\parindent}{0pt}
    \setlength{\parskip}{6pt plus 2pt minus 1pt}}
}{% if KOMA class
  \KOMAoptions{parskip=half}}
\makeatother
\setlength{\emergencystretch}{3em} % prevent overfull lines
\providecommand{\tightlist}{%
  \setlength{\itemsep}{0pt}\setlength{\parskip}{0pt}}
% Slide layout defaults for warning-clean scientific decks.
\usepackage{etoolbox}
\IfFileExists{xurl.sty}{\usepackage{xurl}}{}
\IfFileExists{seqsplit.sty}{\usepackage{seqsplit}}{\newcommand{\seqsplit}[1]{#1}}
\protected\def\breakseq#1{\seqsplit{#1}}
\protected\def\breaktt#1{\begingroup\ttfamily\seqsplit{#1}\endgroup}
\setlength{\emergencystretch}{6em}
\tolerance=5000
\hbadness=10000
\hfuzz=1pt
\setlength{\tabcolsep}{2pt}
\AtBeginEnvironment{longtable}{\tiny\renewcommand{\arraystretch}{0.86}\setlength{\tabcolsep}{1pt}}
\AtBeginEnvironment{tabular}{\tiny\renewcommand{\arraystretch}{0.86}\setlength{\tabcolsep}{1pt}}

% Natbib and cross-reference fallbacks — slides don't load natbib
% or cleveref, but combined-PDF manuscript prose may emit these
% commands. The fallback renders citations as a bracketed key list
% and cross-references as detokenized labels so slides don't fail on
% undefined control sequences or raw underscores. \providecommand is
% a no-op if packages load later.
\providecommand{\citep}[1]{[#1]}
\providecommand{\citet}[1]{#1}
\providecommand{\citealp}[1]{#1}
\providecommand{\citeauthor}[1]{#1}
\providecommand{\citeyear}[1]{#1}
\providecommand{\cref}[1]{\texttt{\detokenize{#1}}}
\providecommand{\Cref}[1]{\texttt{\detokenize{#1}}}
\usepackage{bookmark}
\IfFileExists{xurl.sty}{\usepackage{xurl}}{} % add URL line breaks if available
\urlstyle{same}
% fallback for those not using the hyperref driver hyperxmp:
\makeatletter
\@ifundefined{xmpquote}{\newcommand{\xmpquote}[1]{#1}}{}
\makeatother
\hypersetup{
  hidelinks,
  pdfcreator={LaTeX via pandoc}}

\author{\texorpdfstring{}{}}
\date{}

\begin{document}

\section{Methods Overview}\label{methods-overview}

\begin{frame}[fragile,allowframebreaks]{Methods Overview}
The pipeline is a sequence of deterministic stages, each reading the
previous stage's committed artifacts and writing its own. Business logic
lives in tested \texttt{src/} modules; the numbered \texttt{scripts/}
are thin orchestrators that wire I/O, configuration loading, logging,
and stage sequencing. The architecture follows the thin orchestrator
pattern: no computational logic resides in scripts.
\end{frame}

\begin{frame}[fragile,allowframebreaks]{Pipeline Stages}
\protect\phantomsection\label{pipeline-stages}
\begin{enumerate}
\item
  \textbf{Retrieval} (\breaktt{01\_literature\_search.py}) --- dispatch
  the configured query across 7 engines (arXiv, OpenAlex, Semantic
  Scholar, Crossref, PubMed, SovietRxiv, and ChinaRxiv), merge, and
  de-duplicate into \texttt{corpus.jsonl}. Each engine is an isolated
  adapter exposing a uniform
  \texttt{search(query)\ -\textgreater{}\ list{[}Paper{]}} interface;
  engines that are keyless need no credentials, while Semantic Scholar
  uses a key when present. SovietRxiv and ChinaRxiv share a unified API
  with an optional \texttt{X-API-Email} header for the polite rate-limit
  pool (300/min vs 30/min anonymous).
\item
  \textbf{Meta-analysis} (\breaktt{02\_meta\_analysis\_pipeline.py}) ---
  subfield classification, temporal metrics, TF-IDF, non-negative matrix
  factorization topics, and the citation network. This stage reads
  \texttt{corpus.jsonl} and emits
  \breaktt{subfield\_classification.json},
  \breaktt{temporal\_analysis.json}, \texttt{tfidf\_data.json},
  \texttt{topics.json}, \breaktt{citation\_network.json}, and
  \breaktt{citation\_graph.gml}.
\item
  \textbf{Knowledge graph} (\breaktt{03\_build\_knowledge\_graph.py},
  optional/LLM-gated) --- extract assertions and score the 6 configured
  hypotheses. Outputs \breaktt{nanopublications.jsonl},
  \breaktt{hypothesis\_scores.json}, and
  \breaktt{assertion\_summary.json}.
\item
  \textbf{Figures} (\breaktt{04\_generate\_figures.py}) --- render 18
  publication-ready visualizations from the analysis JSON outputs. All
  figures use a colourblind-safe palette (Wong 2011), high-contrast
  labels at \(\geq 16\)pt, and a headless matplotlib backend (Agg).
\item
  \textbf{Injection} (\breaktt{05\_inject\_variables.py}) --- compute
  manuscript variables from the artifacts above and substitute them into
  these Markdown sections. An unresolved placeholder is a hard error,
  not a silent gap.
\item
  \textbf{Fulltext assessment} (\breaktt{06\_fulltext\_assessment.py})
  --- report abstract coverage (55.5\%), open-access status (14.4\%),
  and PDF availability (40.9\%) across the corpus.
\end{enumerate}
\end{frame}

\begin{frame}[fragile,allowframebreaks]{Reproducibility Model}
\protect\phantomsection\label{reproducibility-model}
The system runs \textbf{offline and deterministically} by default: a
committed synthetic seed corpus drives every stage with fixed seeds
(seed = 42 for NMF, SVD, and graph layouts), so re-running produces
byte-identical outputs. A live run with engines enabled and credentials
supplied replaces the seed corpus with real records --- as in this
instance, which retrieved 2302 live records. The template is
domain-agnostic: the search term, query, keyword set, subfield taxonomy,
and hypotheses all come from \breaktt{manuscript/config.yaml}.
\end{frame}

\begin{frame}[fragile,allowframebreaks]{Configuration Surface}
\protect\phantomsection\label{configuration-surface}
A single \breaktt{manuscript/config.yaml} controls:

\begin{itemize}
\tightlist
\item
  \textbf{Search parameters}: term, query string, per-engine queries,
  relevance keywords, start year, max results, resume/clear behaviour
\item
  \textbf{Engine toggles}: arXiv, OpenAlex, Semantic Scholar, Crossref,
  PubMed, SovietRxiv, ChinaRxiv (each independently enabled or disabled)
\item
  \textbf{SovietRxiv/ChinaRxiv settings}: optional \texttt{api\_email}
  for the polite pool, \texttt{source} filter (\texttt{russiarxiv} or
  \texttt{chinaxiv})
\item
  \textbf{Full-text download}: opt-in Unpaywall resolution with
  \texttt{unpaywall\_email}
\item
  \textbf{Embeddings}: method (\texttt{tfidf\_svd} or
  \texttt{transformer}), dimensionality, max features
\item
  \textbf{Knowledge graph}: checkpoint interval, LLM model, base URL,
  temperature, max tokens
\item
  \textbf{Hypothesis definitions}: 6 named hypotheses with scope labels
\item
  \textbf{Subfield taxonomy}: 6 buckets, each with a keyword list
\item
  \textbf{Paper metadata}: title, authors, DOI, keywords, license,
  repository URL
\end{itemize}
\end{frame}

\end{document}
